\begin{titlepage}
\newgeometry{top=0.3cm}
\newcommand{\HRule}{\rule{\linewidth}{0.5mm}} % Définir une commande pour les lignes horizontales

\center % Centrer tout le contenu de la page

%----------------------------------------------------------------------------------------
%	EN-TÊTE
%----------------------------------------------------------------------------------------

\begin{minipage}[c]{0.18\textwidth}
    \includegraphics[width=\linewidth]{LogoUSo.jpg}
\end{minipage}
\hfill
\begin{minipage}[c]{0.6\textwidth}
    \centering
    \footnotesize
    \textbf{\textsc{République tunisienne}}\\
    \textbf{\textsc{Ministère de l'Enseignement supérieur et de la Recherche scientifique}}\\
    \medskip
    \textbf{\textsc{Université de Sousse}}\\
    \textbf{\textsc{Institut supérieur des sciences appliquées et de technologie de Sousse}}

\end{minipage}
\hfill
\begin{minipage}[c]{0.18\textwidth}
    \includegraphics[width=\linewidth]{ISSAT.png}
\end{minipage}

\vskip1cm
%----------------------------------------------------------------------------------------
%	TITRE DU RAPPORT
%----------------------------------------------------------------------------------------

% TODO: adapter au type de rapport (rapport de stage, de projet, de PFE, etc.).
\textbf{\textsc{\LARGE \bfseries Rapport de stage}}\\[0.5cm]
\vskip0.5cm
% TODO: préciser la classe et l'année d'études.
\textsc{\large Classe : [ANNÉE / CYCLE, par exemple 2e année du cycle ingénieur]}\\
\vskip0.5cm

%----------------------------------------------------------------------------------------
%	AUTEUR ET SUJET
%----------------------------------------------------------------------------------------

\textit{Présenté par :}\\[0.5cm]
\textsc{\large [VOTRE NOM COMPLET]}\\[0.5cm]

\HRule \\[0.4cm]
\textit{\LARGE \bfseries Titre du rapport :}\\[0.4cm]
\textit{\LARGE \bfseries [TITRE DU PROJET / STAGE]}\\
\HRule \\[1cm]

%----------------------------------------------------------------------------------------
%	ENTREPRISE D'ACCUEIL
%----------------------------------------------------------------------------------------

\textsc{\large Entreprise d'accueil :}\\[0.5cm]
% TODO: placer le logo de l'entreprise dans Chapter1/images/ et l'insérer ici,
% ou supprimer cette ligne pour un projet exclusivement universitaire.
% \includegraphics[width=0.25\textwidth]{Chapter1/images/CompanyLogo.png}

\vskip1cm

%----------------------------------------------------------------------------------------
%	ENCADRANTS
%----------------------------------------------------------------------------------------

\begin{minipage}[t]{0.45\textwidth}
    \begin{flushleft} \large
        \emph{Encadrant universitaire :}\\{}
        [NOM DE L'ENCADRANT]
    \end{flushleft}
\end{minipage}
~
\begin{minipage}[t]{0.45\textwidth}
    \begin{flushright} \large
        \emph{Encadrant en entreprise :}\\{}
        [NOM DE L'ENCADRANT]
    \end{flushright}
\end{minipage}

\vfill

% TODO: préciser l'année universitaire.
{\large [ANNÉE UNIVERSITAIRE, par exemple 2026 -- 2027]}

\end{titlepage}
\restoregeometry
